\chapter*{Wykaz skrótów i oznaczeń}
\addcontentsline{toc}{chapter}{Wykaz skrótów i oznaczeń}

% Źródło: wspolne/slownik.md (sekcje 2 i 3), stan 2026-11-09.
% Nowy skrót lub symbol dopisywać najpierw do słownika, potem tutaj.
% Pozycje nieużyte ostatecznie w tekście usunąć przed oddaniem pracy.
% Wymaga: booktabs.

\section*{Skróty}

\begin{tabular}{@{}lp{0.78\textwidth}@{}}
\toprule
\textbf{Skrót} & \textbf{Rozwinięcie} \\
\midrule
APXS    & ang. \textit{Alpha Particle X-ray Spectrometer} -- spektrometr rentgenowski wzbudzany cząstkami alfa \\
CNT     & ang. \textit{Carbon Nanotubes} -- nanorurki węglowe \\
CNT-PVP & kompozyt nanorurek węglowych i~poliwinylopirolidonu \\
DRIE    & ang. \textit{Deep Reactive Ion Etching} -- głębokie trawienie reaktywne jonowe \\
EDS     & ang. \textit{Energy Dispersive X-ray Spectroscopy} -- spektroskopia rentgenowska z~dyspersją energii \\
F--N    & Fowler--Nordheim (równanie, wykres) \\
FWHM    & ang. \textit{Full Width at Half Maximum} -- szerokość połówkowa \\
KOH     & wodorotlenek potasu \\
MCA     & ang. \textit{Multichannel Analyzer} -- analizator wielokanałowy \\
MEMS    & ang. \textit{Micro-Electro-Mechanical Systems} -- mikrosystemy elektromechaniczne \\
NdFeB   & stop neodym--żelazo--bor (magnes trwały) \\
PIXL    & ang. \textit{Planetary Instrument for X-ray Lithochemistry} -- przyrząd łazika Perseverance \\
PVP     & poliwinylopirolidon \\
SDD     & ang. \textit{Silicon Drift Detector} -- krzemowy detektor dryfowy \\
SEM     & ang. \textit{Scanning Electron Microscope} -- elektronowy mikroskop skaningowy \\
Si(Li)  & detektor krzemowy kompensowany litem \\
TMAH    & ang. \textit{Tetramethylammonium Hydroxide} -- wodorotlenek tetrametyloamoniowy \\
UV      & ang. \textit{Ultraviolet} -- nadfiolet \\
WDS     & ang. \textit{Wavelength Dispersive X-ray Spectroscopy} -- spektroskopia z~dyspersją długości fali \\
XRF     & ang. \textit{X-ray Fluorescence} -- fluorescencja rentgenowska \\
ZAF     & poprawki na liczbę atomową (Z), absorpcję (A) i~fluorescencję (F) \\
j.u.    & jednostki umowne \\
\bottomrule
\end{tabular}

\section*{Oznaczenia}

\begin{tabular}{@{}lp{0.62\textwidth}l@{}}
\toprule
\textbf{Symbol} & \textbf{Znaczenie} & \textbf{Jednostka} \\
\midrule
E, B, T                 & emiter, bramka (elektroda ekstrakcyjna), target (elektroda przyspieszająca) & -- \\
$U_\mathrm{B}$          & napięcie ekstrakcyjne (emiter--bramka) & V \\
$U_\mathrm{T}$          & napięcie przyspieszające (emiter--target) & V \\
$I_\mathrm{e}$          & prąd emisji & µA \\
$I_\mathrm{B}$, $I_\mathrm{T}$ & prąd bramki, prąd targetu & µA \\
$R_\mathrm{B}$, $R_\mathrm{T}$ & rezystory ograniczające w~obwodzie bramki i~targetu & k$\Omega$ \\
$L$                     & jasność świecenia ekranu luminescencyjnego & j.u. \\
$p$                     & ciśnienie & Pa \\
$\varphi$               & praca wyjścia & eV \\
$F$                     & natężenie pola elektrycznego na powierzchni emitera & V/m \\
$J$                     & gęstość prądu emisji & A/m$^2$ \\
$\beta$                 & współczynnik wzmocnienia pola & -- \\
$d$                     & odległość emiter--bramka & µm \\
$E_0$                   & energia elektronów wiązki & keV \\
$E$                     & energia fotonu & keV \\
$E_\mathrm{c}$          & energia krytyczna (krawędź absorpcji) & keV \\
$U_0$                   & współczynnik nadnapięcia $E_0/E_\mathrm{c}$ & -- \\
$\vec F_\mathrm{L}$     & siła Lorentza & N \\
$B$                     & indukcja magnetyczna & T \\
$r_\mathrm{L}$          & promień Larmora & mm \\
$f$                     & ogniskowa soczewki magnetycznej & mm \\
$Z$, $A$, $\rho$        & liczba atomowa, masa atomowa, gęstość & --, g/mol, g/cm$^3$ \\
$R_\mathrm{KO}$         & zasięg elektronów wg Kanayi--Okayamy & µm \\
$\omega$                & wydajność fluorescencji & -- \\
$\lambda$               & długość fali & nm \\
$\bar l$                & średnia droga swobodna & mm \\
$\varepsilon$           & średnia energia wytworzenia pary elektron--dziura & eV \\
$\mathcal F$            & czynnik Fano & -- \\
$N$                     & liczba zliczeń & -- \\
$e$, $m_\mathrm{e}$, $c$, $h$, $k_\mathrm{B}$ & ładunek elementarny, masa elektronu, prędkość światła, stała Plancka, stała Boltzmanna & -- \\
\bottomrule
\end{tabular}
